% =============================================================================
% NASA SPACE GRANT PROPOSAL - JACOB VAUGHT
% =============================================================================
%
% COMPILER: pdfLaTeX (recommended)
%
% COMPILE COMMAND:
%   pdflatex vaught_nasa_proposal.tex
%
% =============================================================================

\documentclass[12pt]{article}

% --- Page Layout (1-inch margins required by most grants) ---
\usepackage[margin=1in]{geometry}

% --- Font Encoding & Bold Small Caps ---
\usepackage[T1]{fontenc}
\usepackage{bold-extra}

% --- Colors (UofSC Brand) ---
\usepackage{xcolor}
\definecolor{garnet}{HTML}{73000A}

% --- Paragraph Formatting ---
\setlength{\parindent}{0.25in}
\setlength{\parskip}{0.3em}

% --- Custom Section Commands ---
\newcommand{\sect}[1]{\noindent{\color{garnet}\textbf{\textsc{#1}}.}}
\newcommand{\subsect}[1]{\noindent{\color{garnet}\textit{#1.}}}

% =============================================================================
\begin{document}
{\color{garnet}\large\textbf{\textsc{FPGA-based Digital Twin Modeling of Power Electronics using FINN}}}\\[-0.8em]
{\color{garnet}\rule{\textwidth}{0.5pt}}\\[0.2em]
\textbf{Principal Investigator:} Dr. Austin Downey\\
\textbf{Primary Researcher:} Jacob C. Vaught\\[0.5em]
\sect{Abstract} This project aims to implement FINN on a PYNQ board FPGA for digital twin modeling of power electronics. FINN is an experimental framework from Xilinx Research Labs to explore deep neural network inference on FPGAs. It specifically targets quantized neural networks, with emphasis on generating dataflow-style architectures customized for each network. Power electronics is a field that deals with the conversion and control of electric power using electronic devices and systems. Power electronics has enabled various applications in renewable energy, smart grid, electric vehicles and more. However, power electronic devices and systems are also prone to faults and failures, which can cause serious consequences such as power outage, fire hazard and equipment damage. Therefore, it is essential to monitor and diagnose the health status of power electronic devices and systems in real time and take preventive or corrective actions when needed. Digital twin is an emerging technology that can create a virtual replica of a physical system and synchronize its behavior and performance with the real system using data and feedback. By using digital twin, one can perform various tasks such as simulation, optimization, prediction and control for the physical system without affecting its normal operation. Digital twin can also enable remote monitoring and diagnosis of the physical system by providing a digital interface for data analysis and visualization.

\sect{Specific Objectives} The specific objectives of this project are:
\begin{itemize}
    \setlength{\itemsep}{0.2em}
    \item To train a quantized neural network for a specific power electronic device or system using machine learning techniques such as Brevitas.
    \item To use FINN to generate a dataflow-style FPGA accelerator for the quantized neural network.
    \item To deploy and execute the FPGA accelerator on a PYNQ board using Python drivers and Jupyter notebooks.
    \item To evaluate the performance and accuracy of the FPGA accelerator by comparing it with simulation results or experimental data.
    \item To analyze the advantages and disadvantages of using FINN on PYNQ for digital twin modeling in power electronics.
\end{itemize}

\sect{Expected Goals/Outcomes} The expected goals/outcomes of this project are:
\begin{itemize}
    \setlength{\itemsep}{0.2em}
    \item A novel FPGA-based digital twin technique that can provide real-time monitoring and diagnosis of power electronic devices and systems using quantized neural networks.
    \item A prototype of an FPGA-based digital twin system that can demonstrate the feasibility and effectiveness of the proposed technique.
    \item A comprehensive analysis of the benefits and challenges of using FINN on PYNQ for digital twin modeling in power electronics.
    \item A contribution to the state-of-the-art knowledge and practice in the fields of FPGA, machine learning and digital twin.
\end{itemize}

\sect{NASA Relevance} This project is relevant to NASA's Science mission directorate because it supports NASA's vision to explore the secrets of the universe for the benefit of all; aligns with NASA's mission to explore the unknown in air and space, innovate for the benefit of humanity, and inspire the world through discovery; and provides relevant contributions towards solving NASA Mission Directorate challenges and aligning with the agency's top research priorities. Specifically, this project addresses the following NASA research areas:
\begin{itemize}
    \setlength{\itemsep}{0.2em}
    \item Heliophysics: Understanding how power electronics can be used to improve space weather forecasting and mitigation.
    \item Planetary Science: Developing robust and reliable power electronics for planetary exploration missions.
    \item Astrophysics: Enhancing power electronics for advanced telescopes and observatories.
    \item Earth Science: Improving power electronics for renewable energy systems and climate monitoring.
\end{itemize}

\sect{Methods} This research follows a five-step process. First, we train a quantized neural network for a specific power electronic device using Brevitas. Second, we generate a dataflow-style FPGA accelerator for the network with FINN, optimizing resources and performance.

In the third step, we deploy the accelerator on a PYNQ board using Python drivers and Jupyter notebooks. Fourth, we evaluate the FPGA accelerator's performance and accuracy by comparing it to simulation results or experimental data, assessing throughput, latency, and power consumption.

Finally, we analyze the pros and cons of using FINN on PYNQ for digital twin modeling in power electronics, comparing it to existing techniques, addressing challenges, and suggesting improvements for future work.

\end{document}
